% Shared clean-flat-vector SysML style for the survey model diagrams.
% Included by each standalone diagram source; compiled to a vector PDF that
% replaces the original Cameo raster export. No drop shadows, no gradients:
% solid peach header, cream body, thin border.
\usepackage{tikz}
\usetikzlibrary{shapes.multipart,positioning,fit,backgrounds,calc,arrows.meta}

\definecolor{cameoHeader}{RGB}{247,206,151}
\definecolor{cameoBody}{RGB}{253,246,225}
\definecolor{cameoBorder}{RGB}{170,120,70}
\definecolor{cameoLit}{RGB}{120,120,120}
\definecolor{frameBorder}{RGB}{90,90,90}
\definecolor{reqFill}{RGB}{245,220,235}
\definecolor{reqBorder}{RGB}{150,90,140}
\definecolor{pkgFill}{RGB}{228,228,247}
\definecolor{pkgBorder}{RGB}{120,120,175}

% UML package container: call AFTER its inner nodes exist.
% #1 = fit list e.g. (a)(b)(c), #2 = container node id, #3 = package title
\newcommand{\pkgframe}[3]{%
  \begin{scope}[on background layer]
    \node[draw=pkgBorder,line width=0.6pt,fill=pkgFill,fit=#1,
          inner xsep=5mm,inner ysep=8mm] (#2) {};
  \end{scope}
  % title near the top and a folder tab on the top-left corner
  \node[anchor=north,font=\sffamily\bfseries] at ([yshift=-1.5mm]#2.north) {#3};%
  \node[draw=pkgBorder,line width=0.6pt,fill=pkgFill,anchor=south west,
        minimum width=10mm,minimum height=2.8mm,inner sep=0pt]
        at ([xshift=0.3pt,yshift=-0.3pt]#2.north west) {};%
}

% small inline «requirement» marker, e.g. Satisfies = \reqmark\ Research Question 1
\newcommand{\reqmark}{{\setlength{\fboxsep}{1pt}\setlength{\fboxrule}{0.4pt}%
  \fcolorbox{reqBorder}{reqFill}{\tiny\bfseries R}}}

\tikzset{
  sysmlbase/.style={
    rectangle split, draw=cameoBorder, line width=0.6pt,
    rectangle split draw splits=true,
    inner xsep=8pt, inner ysep=6pt,
    font=\sffamily, align=center, anchor=north west,
  },
  enumnode/.style={
    sysmlbase, rectangle split parts=2,
    rectangle split part fill={cameoHeader,cameoBody},
    text width=36mm,
  },
  blocknode/.style={
    sysmlbase, rectangle split parts=2,
    rectangle split part fill={cameoHeader,cameoBody},
    text width=52mm,
  },
  blocknodewide/.style={
    sysmlbase, rectangle split parts=3,
    rectangle split part fill={cameoHeader,cameoBody,cameoBody},
    text width=150mm,
  },
  taxnode/.style={
    rectangle, draw=cameoBorder, line width=0.6pt, fill=cameoHeader,
    inner xsep=8pt, inner ysep=5pt, font=\sffamily, align=center,
    anchor=north west, text width=44mm,
  },
  reqnode/.style={
    sysmlbase, rectangle split parts=2, draw=reqBorder,
    rectangle split part fill={reqFill,reqFill!55!white},
    text width=58mm,
  },
}
% compact «Block» node for the (dense) requirement-traceability diagram
\tikzset{
  cblocknode/.style={
    rectangle, draw=cameoBorder, line width=0.5pt, fill=cameoHeader,
    inner xsep=4pt, inner ysep=3pt, font=\sffamily, align=center,
    anchor=north west, text width=20mm,
  },
  satisfyedge/.style={dashed, draw=reqBorder, line width=0.5pt, -{Triangle[open,length=6pt,width=7pt]}},
}
\newcommand{\cblock}[3]{\node[cblocknode,#1] (#2) {\tiny\guillemotleft Block\guillemotright\\[0.3pt]\scriptsize\bfseries #3};}

% single-compartment «block» taxonomy node: #1 opts, #2 id, #3 name
\newcommand{\taxbox}[3]{\node[taxnode,#1] (#2) {\footnotesize\guillemotleft block\guillemotright\\[1pt]\bfseries #3};}
% «requirement» box: #1 opts, #2 id, #3 name, #4 body (Id/Text, left-aligned)
\newcommand{\reqbox}[4]{%
  \node[reqnode,#1] (#2) {%
    \nodepart[align=center]{one}{\footnotesize\guillemotleft requirement\guillemotright\\[1pt]\bfseries #3}%
    \nodepart[align=left]{two}{\upshape\footnotesize #4}%
  };
}

% compartment label (grey italic, centred): e.g. "enumeration literals", "operations"
\newcommand{\complabel}[1]{{\itshape\scriptsize\textcolor{cameoLit}{\mbox{}\hfill #1\hfill\mbox{}}}}

% «enumeration» box: #1 placement options (e.g. right=6mm of a.north east,anchor=north west),
% #2 node id, #3 name, #4 literals (\\-separated)
\newcommand{\enumbox}[4]{%
  \node[enumnode,#1] (#2) {%
    \nodepart[align=center]{one}{\footnotesize\guillemotleft enumeration\guillemotright\\[1pt]\bfseries #3}%
    \nodepart[align=left]{two}{\complabel{enumeration literals}\\[2pt]\upshape\normalsize #4}%
  };
}

% «Block» box (2 compartments): #1 placement options, #2 id, #3 name, #4 operations
\newcommand{\blockbox}[4]{%
  \node[blocknode,#1] (#2) {%
    \nodepart[align=center]{one}{\footnotesize\guillemotleft Block\guillemotright\\[1pt]\bfseries #3}%
    \nodepart[align=left]{two}{\complabel{operations}\\[2pt]\upshape\normalsize #4}%
  };
}

% «Block» box (3 compartments, detailed): #1 placement options, #2 id, #3 name, #4 properties, #5 operations
\newcommand{\blockboxwide}[5]{%
  \node[blocknodewide,#1] (#2) {%
    \nodepart[align=center]{one}{\footnotesize\guillemotleft Block\guillemotright\\[1pt]\bfseries #3}%
    \nodepart[align=left]{two}{\upshape\normalsize #4}%
    \nodepart[align=left]{three}{\complabel{operations}\\[2pt]\upshape\normalsize #5}%
  };
}

% Package-diagram frame + title tab. Call AFTER a node named (frame) has been
% placed via \node[fit=...]{}; pass the title text.
\newcommand{\packageframe}[1]{%
  \draw[frameBorder,line width=0.8pt] (frame.south west) rectangle (frame.north east);%
  \node[anchor=north west,frameBorder,draw,line width=0.8pt,fill=white,inner sep=3pt,
        align=left,font=\sffamily] at (frame.north west) {#1};%
}
